%! Skills Focus: Selecting an Appropriate Inference Procedure for Categorical Data — Unit 8, Lesson 8.7 — Warm-Up (Answer Key)
\documentclass[10pt]{article}
\usepackage{apstats-article}
\usepackage{apstats-key}   % pulls in -boxes; do NOT also load -boxes

\begin{document}

\pageheader{Unit 8, Lesson 8.7}{Warm-Up --- Answer Key}
\namedateperiod

\vspace{0.12in}

\begin{notesbox}{Spiral Review --- Identifying Categorical Inference Procedures (Units 6 \& 8)}

\textit{For each scenario, identify the correct inference procedure. No computation needed.}

\vspace{0.06in}
\textbf{1.} A botanist randomly selects 120 seedlings and records whether each sprouted
(yes or no) under a new fertilizer treatment. She claims that more than 60\% of seedlings
will sprout. She wants to test this claim.

\begin{enumerate}[label=\alph*.]
  \item Does the outcome have two categories or more than two? \ans{two (yes/no)}
  \item How many groups or samples? \ans{one}
  \item Is this a test or a confidence interval? \ans{test}
  \item Name the procedure: \ans{one-sample $z$-test for $p$}
\end{enumerate}

\vspace{0.10in}
\textbf{2.} A school district randomly surveys students at three different high schools and
asks each student to pick their favorite core subject (Math, English, Science, or Social
Studies). A researcher wants to test whether the distribution of favorite subjects is the
same at all three schools.

\begin{enumerate}[label=\alph*.]
  \item How many categories does the outcome variable have? \ans{four}
  \item How many independent groups are being compared? \ans{three}
  \item Name the procedure: \ans{chi-square test for homogeneity}
\end{enumerate}

\vspace{0.10in}
\textbf{3.} A quality-control manager randomly selects 200 cereal boxes from a factory line and
records whether each box is underfilled, correctly filled, or overfilled. The manufacturer
claims the boxes should be 5\% underfilled, 90\% correctly filled, and 5\% overfilled. She
wants to test whether the actual distribution matches this claim.

Name the procedure: \ans{chi-square goodness-of-fit test}

\end{notesbox}

\end{document}
